\documentclass[12pt,a4paper]{article}

% ── Packages ──────────────────────────────────────────────────────────────────
\usepackage[utf8]{inputenc}
\usepackage[T1]{fontenc}
\usepackage{amsmath,amssymb}
\usepackage{graphicx}
\usepackage{booktabs}
\usepackage{hyperref}
\usepackage{natbib}
\usepackage{geometry}
\usepackage{float}
\usepackage{caption}
\usepackage{subcaption}
\usepackage{algorithm}
\usepackage{algpseudocode}
\usepackage{xcolor}
\geometry{margin=1in}

% ── Metadata ──────────────────────────────────────────────────────────────────
\title{%
  \textbf{ColdGuard: A Kinetic-Bayesian Decision Support System for \\
  Vaccine Potency Estimation Under Cold Chain Excursions}%
}

\author{
  [Author Name] \\
  Department of [Department], [Institution] \\
  \texttt{[email@institution.ac.in]}
}

\date{[Month Year]}

% ── Document ──────────────────────────────────────────────────────────────────
\begin{document}

\maketitle

\begin{abstract}
Vaccine wastage due to temperature excursions in cold chains is a critical
challenge for immunisation programmes in low- and middle-income countries.
Current protocols---based on Vaccine Vial Monitors (VVMs) or simple threshold
rules---produce high False Discard Rates (FDR), destroying viable vaccine, or
high False Use Rates (FUR), risking under-potent doses. We present
\textbf{ColdGuard}, a kinetic-Bayesian decision support system that integrates
Arrhenius degradation kinetics with Monte Carlo uncertainty propagation to
produce calibrated potency estimates and confidence-qualified USE/INVES\-TIGATE/DISCARD
decisions from arbitrary temperature logs. Evaluated on [N] simulated cold chain
scenarios representative of India's district, PHC, and outreach tiers, ColdGuard
achieves an FDR of [X\%] and FUR of [Y\%] compared to [A\%] FDR and [B\%] FUR
for WHO Mean Kinetic Temperature (MKT) and [C\%] FDR and [D\%] FUR for VVM-only
protocols, while providing interpretable natural-language explanations suitable
for use by Auxiliary Nurse Midwives with no specialist training. An offline
mobile application for Android implements the same engine at 1000 Monte Carlo
samples, running to completion within 3 seconds on a low-cost device. ColdGuard
is open-source under MIT licence.
\end{abstract}

\textbf{Keywords:} vaccine potency; cold chain; Arrhenius kinetics; Bayesian
uncertainty quantification; decision support; immunisation; India; last-mile

\tableofcontents
\newpage

% ══════════════════════════════════════════════════════════════════════════════
\section{Introduction}
\label{sec:introduction}

[TODO: Motivate the problem. Vaccine-preventable disease burden. India's UIP
scale (390 million doses/year). Cold chain failures at last-mile. Cost of FDR
(vaccine waste) and FUR (programme trust, effectiveness). State the gap:
existing tools are binary threshold rules with no uncertainty quantification.
State the contribution: ColdGuard + evaluation.]

% ── Background ────────────────────────────────────────────────────────────────
\section{Background}
\label{sec:background}

\subsection{Arrhenius Kinetics for Pharmaceutical Degradation}

Vaccine degradation follows first-order Arrhenius kinetics \citep{galazka1998thermostability}.
For a time-varying temperature trajectory $T(t)$, the remaining potency fraction is:

\begin{equation}
  P(t) = P_0 \cdot \exp\!\left(-\int_0^t k\bigl(T(\tau)\bigr)\,d\tau\right),
  \quad k(T) = A \cdot \exp\!\!\left(\frac{-E_a}{RT}\right)
  \label{eq:arrhenius-potency}
\end{equation}

where $E_a$ [J/mol] is the activation energy, $A$ [hr$^{-1}$] is the
pre-exponential factor, and $R = 8.314$ J/(mol$\cdot$K).

For a discrete temperature log $\{(t_i, T_i)\}$, we evaluate the integral using
the trapezoidal rule:

\begin{equation}
  \int_0^t k(T(\tau))\,d\tau \approx \sum_{i=1}^{n-1}
  \frac{k(T_i) + k(T_{i+1})}{2}\,\Delta t_i
  \label{eq:trapezoid}
\end{equation}

\subsection{WHO Mean Kinetic Temperature}

The WHO MKT \citep{who2006temperature} gives a single equivalent temperature:

\begin{equation}
  T_\text{MK} = \frac{E_a/R}{-\ln\!\left[\frac{1}{n}\sum_{i=1}^n
    \exp\!\!\left(\frac{-E_a}{RT_i}\right)\right]}
  \label{eq:mkt}
\end{equation}

MKT provides a baseline but yields a single point estimate with no calibrated
uncertainty quantification.

\subsection{Vaccine Vial Monitors}

[TODO: Brief description of VVMs. Their conservative calibration. Trade-off
between sensitivity and specificity. Published FDR estimates from literature.]

% ── Methods ───────────────────────────────────────────────────────────────────
\section{Methods}
\label{sec:methods}

\subsection{System Architecture}

[TODO: Figure~\ref{fig:architecture}. Four-layer pipeline: parsing → Arrhenius
integration → Monte Carlo uncertainty propagation → decision layer.]

\begin{figure}[H]
  \centering
  % \includegraphics[width=0.85\textwidth]{figures/fig1_architecture.pdf}
  \fbox{\parbox{0.85\textwidth}{\centering [PLACEHOLDER: System architecture diagram]}}
  \caption{ColdGuard system architecture. Temperature logs are parsed, degradation
    is integrated via Arrhenius kinetics, uncertainties in $E_a$, $A$, and sensor
    readings are propagated via Monte Carlo sampling, and a calibrated decision is
    emitted.}
  \label{fig:architecture}
\end{figure}

\subsection{Bayesian Uncertainty Propagation}

Arrhenius parameters are uncertain. ColdGuard propagates three uncertainty
sources via Monte Carlo ($N = 5000$ samples):

\begin{enumerate}
  \item \textbf{Activation energy:} $E_a^{(s)} \sim \mathcal{N}(\mu_{E_a}, \sigma_{E_a}^2)$
  \item \textbf{Pre-exponential factor:} $\ln A^{(s)} \sim \mathcal{N}(\ln A_\text{mean}, \sigma_{\ln A}^2)$
  \item \textbf{Sensor error:} $\varepsilon_i^{(s)} \sim \mathcal{N}(0,\,\sigma_\text{logger}^2)$, added to each reading
\end{enumerate}

The posterior potency distribution $\{P^{(1)}, \ldots, P^{(N)}\}$ yields the
90\% credible interval and the key probability $\Pr(P > P_\text{min})$ that drives
the decision.

\subsection{Decision Logic}

\begin{table}[H]
  \centering
  \caption{Decision thresholds}
  \begin{tabular}{lll}
    \toprule
    Condition & Decision \\
    \midrule
    $\Pr(P > P_\text{min}) > 0.90$ & \textsc{use} \\
    $0.70 < \Pr(P > P_\text{min}) \leq 0.90$ & \textsc{investigate} \\
    $\Pr(P > P_\text{min}) \leq 0.70$ & \textsc{discard} \\
    \bottomrule
  \end{tabular}
  \label{tab:decision}
\end{table}

\subsection{Vaccine Parameters}

\begin{table}[H]
  \centering
  \caption{Arrhenius parameters for UIP vaccines (sources: \citealt{galazka1998thermostability,who2006temperature})}
  \begin{tabular}{lrrrr}
    \toprule
    Vaccine & $E_a$ (kJ/mol) & Shelf life & Ref.\ temp.\ (°C) & Min.\ potency \\
    \midrule
    DPT       & $83 \pm 4$   & 24 months & 5  & 80\% \\
    OPV       & $112 \pm 6$  & 6 months  & 5  & 67\% \\
    MMR       & $108 \pm 7$  & 12 months & 5  & 80\% \\
    BCG       & $90 \pm 5$   & 12 months & 5  & 80\% \\
    Hep B     & $70 \pm 4.5$ & 24 months & 5  & 80\% \\
    IPV       & $92 \pm 5.5$ & 24 months & 5  & 80\% \\
    Rotavirus & $100 \pm 6$  & 24 months & $-15$ & 80\% \\
    PCV       & $77 \pm 4$   & 24 months & 5  & 80\% \\
    \bottomrule
  \end{tabular}
  \label{tab:params}
\end{table}

\subsection{Simulation Study Design}

[TODO: Describe the 3-tier India cold chain simulation (district hospital, PHC,
outreach). Scenario generation parameters. 1000 simulated cold chains. Ground
truth labelling methodology. McNemar's test for comparing ColdGuard vs MKT vs
VVM. FDR/FUR metrics definition.]

\subsection{Validation Against Published Data}

[TODO: Literature validation. Describe sources. Round-trip test methodology.
Calibration analysis (coverage probability, RMSE).]

% ── Evaluation ────────────────────────────────────────────────────────────────
\section{Evaluation}
\label{sec:evaluation}

[TODO: Computational environment. Test suite (37 unit/integration tests).
Runtime benchmarks (mobile: 1000 samples < 3 s). Parameter sensitivity analysis.]

% ── Results ───────────────────────────────────────────────────────────────────
\section{Results}
\label{sec:results}

\subsection{Example Scenario}

[TODO: Figure~\ref{fig:example}: DPT power-outage scenario. Temperature timeline,
potency distribution, segment attribution.]

\begin{figure}[H]
  \centering
  \fbox{\parbox{0.85\textwidth}{\centering [PLACEHOLDER: Example DPT scenario figure]}}
  \caption{ColdGuard output for the demo DPT power-outage scenario (11 h at
    14.3°C). Left: temperature timeline with excursion highlighted. Right:
    Monte Carlo potency distribution with 90\% CI band.}
  \label{fig:example}
\end{figure}

\subsection{Simulation Study Results}

[TODO: Table comparing FDR/FUR/Accuracy for ColdGuard, MKT, VVM across three
cold chain tiers. Figure~\ref{fig:simulation}.]

\begin{figure}[H]
  \centering
  \fbox{\parbox{0.85\textwidth}{\centering [PLACEHOLDER: Simulation results figure]}}
  \caption{FDR and FUR for three decision protocols across India cold chain tiers.}
  \label{fig:simulation}
\end{figure}

\subsection{Calibration}

[TODO: Calibration plot (reliability diagram). Coverage probability table.
RMSE and MAE on validation data.]

% ── Discussion ────────────────────────────────────────────────────────────────
\section{Discussion}
\label{sec:discussion}

[TODO: Interpret results. Why ColdGuard reduces FDR: calibrated uncertainty
avoids conservative thresholding. Why INVESTIGATE zone is valuable vs binary.
Limitations: Ea values from literature starting points, not independently derived;
validation data partially synthetic; mobile not field-tested. Comparison to
prior work: Matthias et al. potency estimation, ColdTrace logger software,
WHO guidance evolution. Deployment considerations: offline-first, Hindi i18n,
low-cost Android.]

\subsection{Limitations}

\begin{enumerate}
  \item \textit{Parameter uncertainty}: Ea values are drawn from published ranges,
    not independently derived from primary stability data for each product.
  \item \textit{Validation data}: Only three vaccines have literature-sourced
    stability points; remaining five use model-generated surrogate data.
  \item \textit{Freeze damage}: The current model does not quantify damage from
    freeze events; only detection is implemented.
  \item \textit{Field validation}: The mobile application has not been tested
    in a real cold chain setting with ANM users.
\end{enumerate}

% ── Conclusion ────────────────────────────────────────────────────────────────
\section{Conclusion}
\label{sec:conclusion}

[TODO: Summary. ColdGuard demonstrates that calibrated Bayesian uncertainty
quantification applied to first-principles Arrhenius kinetics meaningfully
reduces both false discard and false use rates compared to VVM and MKT protocols,
while remaining deployable offline on low-cost Android devices. Future work:
field validation, real stability data for all 8 vaccines, freeze damage model,
integration with national eVIN cold chain monitoring system.]

% ── References ────────────────────────────────────────────────────────────────
\bibliographystyle{plainnat}
\bibliography{references}

% ── Appendix ──────────────────────────────────────────────────────────────────
\appendix
\section{Supplementary Material}
\label{sec:supplementary}

See the \texttt{paper/supplementary/} directory for:
\begin{itemize}
  \item \texttt{S1\_arrhenius\_derivation.pdf} — Full derivation of A from shelf-life specification
  \item \texttt{S2\_mc\_convergence.pdf} — Monte Carlo convergence diagnostics
  \item \texttt{S3\_scenario\_descriptions.pdf} — All 10 synthetic test scenario descriptions
\end{itemize}

\end{document}
